\documentclass[10pt,a4paper]{amsart}
\usepackage[margin=19mm]{geometry}
\usepackage{amsmath,amssymb,mathtools}
\usepackage[scr=rsfs,cal=euler]{mathalfa}
\usepackage{xfrac}
\usepackage[unicode,hidelinks,bookmarks=false]{hyperref}
\newcommand{\ie}{i.e.\ }
\newcommand{\cf}{cf.\ }
\newcommand{\resp}{respectively\ }
\newcommand{\Subsection}{\subsection}
\providecommand{\nobreakdash}{\nobreak\hbox{-}}
\setlength{\emergencystretch}{2em}
\multlinegap0pt
\usepackage{amssymb}
\usepackage{xcolor}
\newtheorem{prop}{Proposition}
\newtheorem{lema}{Lemma}
\DeclareMathOperator{\D}{D^2}
\DeclareMathOperator{\I}{I}
\DeclareMathOperator{\LP}{Conv_{lp}}
\DeclareMathOperator{\LQ}{_{l.q}}
\newcommand{\RR}{(\mathbb{R}^n)}
\DeclareMathOperator{\diam}{diam}
\title{A Characterization of Functional Affine Surface Areas}
\author{Fernanda M. Baêta}
\date{Source: arXiv:2512.08375v2 (CC BY 4.0)}
\begin{document}
\maketitle

\noindent\emph{Source and licence.} A Characterization of Functional Affine Surface Areas. Version arXiv:2512.08375v2 (CC BY 4.0), distributed by arXiv under the Creative Commons Attribution 4.0 International licence, http://creativecommons.org/licenses/by/4.0/, as recorded in the arXiv licence metadata for this exact version and shown on the abstract page https://arxiv.org/abs/2512.08375v2. Full licence notice: \texttt{licenses/arxiv-2512.08375-CC-BY-4.0.txt}.\par\smallskip

\noindent\textit{Editorial scope.} Counted unit: the article's prop regularr (source0.tex lines 2027--2047). The article's complete original proof of it is delivered with it (source0.tex lines 2381--2416, one \texttt{proof} environment of 36 lines, which the article places away from the statement). No mathematics is written, repaired or re-derived: the statement and the proof are the article's own lines, in the article's own notation, and no retained line is substituted or re-flowed.  Retained uncounted as the source-local context the proof needs: lines 1559--1561; lines 2090--2092; lines 2159--2163; lines 2301--2310; lines 2326--2335; lines 2341--2343.  Nothing was omitted from any retained span, so the package carries no omission record.  No retained line contains a citation, so the wrapper carries no bibliography and no bibliography entry.  No retained line contains a figure environment, an image-inclusion command or drawing code, so no image is delivered and the package declares no asset dependency.  Editorial formatting only, in addition: an AMS \texttt{amsart} preamble that restates the article's own declarations and macros used by the retained lines, namely the environments \texttt{prop}, \texttt{lema} on separate counters, together with the standard packages those lines need.  The retained environments print in this document as Proposition 1, Lema 1 in order of appearance, whereas the article prints the same items with the numbering of its own full document.  The complete native source of the article as distributed in the arXiv e-print for this exact version is bundled unchanged as \texttt{sources/190652/source0.tex}.
\begin{align}\label{li_lc}
l^c(x)< u(x)< l^i(x)
\end{align}

\begin{align}\label{eq10.26.}
   \diam(C)^n= n^{\frac{n}{2}}\, V_n(C).
\end{align}

\begin{lema}\label{lema10.1}
Under the above hypotheses,
$$Z(u + \I_{rC})\leq Z(q_u+\I_{rC})+\frac{\rho}{4}V_n(rC)$$
for all sufficiently small $r > 0$.
\end{lema}

\begin{lema}\label{claim2}
For every $r$, $0 < r < 1$, define
$$
G_r = \{ y \in \mathbb{R}^n \mid \mathrm{ext}_{(1-4\sqrt{t}) rC}(q_u)(y) = u (y) \}.
$$
Then
\begin{align*}
G_r \subset rC.
\end{align*}    
\end{lema}

\begin{lema}\label{lema*}
For sufficiently small $r>0$, there exists  $v_r\in P_{\LQ}\RR$ such that
\begin{align*}
 (u + \I_{rC})(x)\leq (v_r+\I_{rC})(x) \leq  (q_u+\I_{rC})(x) 
\end{align*}
and 
\begin{align}\label{eq10.39}
    Z(u + \I_{rC})\leq Z(v_r+\I_{rC})+\dfrac{\rho}{2}V_n(rC).
\end{align}
\end{lema}

\begin{align}\label{eq10.42}
v_r = \mathrm{ext}_{(1-4\sqrt{t}) rC}(q_u) \vee  l^c, 
\end{align}

\begin{prop}\label{regularr}
Let $Z:\LP\RR\rightarrow\mathbb{R}$ be a simple, $\tau$-upper semicontinuous, equi-affine and  dually epi-translation   invariant valuation that   vanishes on  indicator functions of polytopes and cylinder functions. Let $u \in \LP(\mathbb{R}^n)$,
let $x_0\in N$ be such that $\det\D u(x_0)>0$,  let $\rho>0$, and let $l^i, l^c$  be the functions satisfying \eqref{li_lc}. Then there exist $r(x_0)>0$, 
polytopes $P_r(x_0)$ containing $x_0$ in their interiors, and  functions $v_r^{x_0}\in P\LQ\RR$ such that
\begin{align}\label{eq10.47a}
  Z(u+\I_{P_r(x_0)}) \leq Z(v_r^{x_0}+\I_{P_r(x_0)}) + \frac{\rho}{2}V_n(P_r(x_0))
\end{align}
and
\begin{align}\label{eq10.48b}
l^c(x)\leq v_r^{x_0}(x)\leq l^i(x)    
\end{align}
for all $x\in \mathbb{R}^n$ and $0<r\leq r(x_0)$. Moreover, for $0<r\leq r(x_0)$,  the function $v_r^{x_0}$ can be chosen so such that
$u$ and $v_r^{x_0}$ share a common supporting hyperplane at $x_0$, and 
\begin{align}\label{eq4111c}
 \{x\in \mathbb{R}^n\mid v_r^{x_0}(x)\geq u(x)\}   \subset P_r(x_0), 
\end{align}
while the  polytopes $P_r(x_0)$ satisfy 
\begin{align}\label{eq10.26..d}
\diam(P_r(x_0))^n= n^{\frac{n}{2}}\, V_n(P_r(x_0)).    
\end{align}
\end{prop}

\begin{proof}[Proof of Proposition~\ref{regularr}]
Let $r>0$ be sufficiently small so that the conclusions of Lemma~\ref{lema10.1} and Lemma~\ref{lema*} hold. Define
\begin{align}\label{eq10.44}
    P_r(x_0)=rC
    \quad \text{and} \quad
    v_r^{x_0}=v_r.
\end{align}
By Lemma~\ref{lema*}, we have $v_r^{x_0}\in P\LQ\RR$ and
\[
Z(u+\I_{P_r(x_0)})
\le
Z(v_r^{x_0}+\I_{P_r(x_0)})
+\frac{\rho}{2}V_n(P_r(x_0)).
\]
This gives \eqref{eq10.47a}.

Moreover, by the construction of $v_r$ in \eqref{eq10.42}, and by choosing $r>0$ sufficiently small, we have
\[
l^c(x)\leq v_r^{x_0}(x)\leq l^i(x)
\quad \text{for all } x\in\mathbb{R}^n,
\]
which gives \eqref{eq10.48b}. Also, since $v_r$ is defined through the tangential extension of $q_u$ on $(1-4\sqrt{t})rC$ and shares the same supporting hyperplane as $u$ at $x_0$, the functions $u$ and $v_r^{x_0}$ share a common supporting hyperplane at $x_0$. By Lemma~\ref{claim2}, the set where this extension can coincide with $u$ is contained in $rC=P_r(x_0)$; hence
\[
\{x\in\mathbb{R}^n: v_r^{x_0}(x)\geq u(x)\}
\subset P_r(x_0),
\]
which gives \eqref{eq4111c}.

Finally, by \eqref{eq10.26.}, since $P_r(x_0)=rC$, we have
\[
\diam(P_r(x_0))^n
=
n^{\frac n2}V_n(P_r(x_0)).
\]
Thus there exists $r(x_0)>0$ such that \eqref{eq10.47a}, \eqref{eq10.48b}, \eqref{eq4111c}, and \eqref{eq10.26..d} hold for every $0<r\le r(x_0)$. This completes the proof.
\end{proof}

\end{document}
